\documentclass[11pt]{article}
\usepackage[margin=0.88in]{geometry}
\usepackage{amsmath,amssymb,lmodern}
\usepackage[T1]{fontenc}
\usepackage{needspace}
\usepackage[colorlinks=true,linkcolor=blue,urlcolor=blue]{hyperref}
\setlength{\emergencystretch}{3em}
\setlength{\parskip}{0.35em}
\setcounter{secnumdepth}{0}
\providecommand{\tightlist}{\setlength{\itemsep}{0pt}\setlength{\parskip}{0pt}}
\title{A genus obstruction for commuting boundary operators}
\author{Henry Zweiman}
\date{September 15, 2026}
\begin{document}
\maketitle
\section{Abstract}\label{abstract}

Let \(M\) be a compact connected oriented smooth Riemannian surface with
nonempty smooth boundary. We prove that if its Dirichlet-to-Neumann map
commutes with the boundary Laplacian, then \(M\) has genus zero. This
answers both cases of Speciel\textquotesingle s Open Problem 1.6 and,
together with his genus-zero classification, removes the topological
restriction from the classification of oriented surfaces with commuting
boundary operators. For connected boundary, the finite-dimensional
obstruction to harmonic conjugation forces an exact determinant identity
on all sufficiently high Fourier planes. The algebra of holomorphic
boundary traces then produces a holomorphic map of degree one. For
several boundary components, we prove that conformally attaching a
Euclidean disk along an arclength identification preserves commutation.
The proof separates the finite-period input, the multiplicative trace
argument, and the variational elimination of the capped boundary.

\textbf{2020 Mathematics Subject Classification:} 58J50, 35P05, 30F15.

\textbf{Keywords:} Dirichlet-to-Neumann map, boundary Laplacian, genus,
harmonic conjugation, conformal welding.

\section{1. Introduction}\label{introduction}

Let \((M,g)\) be a compact connected smooth Riemannian surface with
nonempty smooth boundary. Its Dirichlet-to-Neumann map is

\[
Af=\partial_\nu u_f\big|_{\partial M},\qquad
\Delta_g u_f=0,\qquad u_f\big|_{\partial M}=f,
\tag{1.1}
\]

where \(\nu\) is the outward unit normal and \(\Delta_g\) is the
nonnegative Laplacian. Write \(\Delta_\partial\) for the nonnegative
Laplacian of the induced boundary metric. We consider the exact
commutation condition

\[
A\Delta_\partial f=\Delta_\partial Af
\qquad\text{for every }f\in C^\infty(\partial M).
\tag{1.2}
\]

Speciel {[}S, Theorem 1.3{]} classified the surfaces satisfying (1.2)
under the additional assumption that the genus is zero or that there are
at least three boundary components. His classification contains disks,
logarithmic ovals, and flat cylinders, with the equivalence appropriate
to two-dimensional Dirichlet-to-Neumann data. The remaining cases, a
positive-genus oriented surface with one or two boundary components, are
explicitly posed in {[}S, Open Problem 1.6{]}. The numbering here refers
to the final journal article; the corresponding question is Open Problem
1.12 in its arXiv version.

\Needspace{6\baselineskip}\textbf{Theorem 1.1 (Genus obstruction).} Let \((M,g)\) be a compact
connected oriented smooth Riemannian surface with nonempty smooth
boundary. If (1.2) holds, then \(M\) has genus zero.

The theorem imposes no curvature condition and no restriction on the
number or lengths of the boundary components. In particular, neither of
the positive-genus surfaces in {[}S, Open Problem 1.6{]} admits the
requested metric.

Recall that two surfaces are \emph{\(\sigma\)-isometric} if a
diffeomorphism identifies their metrics after a conformal change
\(g\mapsto e^{2\varphi}g\) with \(\varphi=0\) on the boundary. Both
operators in (1.2) are preserved by this equivalence. A logarithmic oval
is, up to Euclidean similarities, the image of the closed unit disk
under \(z\mapsto\log(z-a)\), \(a>1\), using a branch holomorphic near
the disk. We use the flat cylinder models of {[}S{]}.

\Needspace{6\baselineskip}\textbf{Corollary 1.2 (Classification without a genus restriction).} A
compact connected oriented smooth Riemannian surface with nonempty
smooth boundary satisfies (1.2) if and only if it is
\(\sigma\)-isometric to a disk, a logarithmic oval, or a flat cylinder.

\emph{Proof.} Theorem 1.1 reduces the assertion to the genus-zero case
of {[}S, Theorem 1.3{]}, which supplies both the classification and its
converse. \(\square\)

The new assertion is the genus obstruction. The models, their
commutation property, and their classification in genus zero belong to
Speciel. Our proof does not start from a uniformization of a
positive-genus surface by a planar domain. Instead, it constructs a
degree-one holomorphic map from the boundary operator identity.

When the boundary is connected, let \(T\) denote positively oriented
arclength differentiation. Harmonic conjugation can fail globally
because of periods, but these periods lie in a finite-dimensional
cohomology space. Consequently, on mean-zero functions,
\((A^{-1}T)^2+I\) has finite rank. This is a classical feature of the
boundary Hilbert transform in the inverse problem for Riemann surfaces;
see {[}B{]} and the formulation in {[}BK, Section 2{]}. We give the
needed period argument directly. Under (1.2), the finite-rank defect
becomes an exact determinant constraint on every sufficiently high
two-dimensional Fourier plane.

The further input is multiplicative. Smooth holomorphic boundary traces
form an algebra, characterized by \(Af=-iTf\). Condition (1.2) makes
that algebra invariant under Fourier-plane projections. Multiplying
adjacent high-frequency traces and projecting to frequency one
constructs an ellipse-valued boundary parametrization. The exceptional
case in which all negative high-frequency coefficients vanish produces a
disk-valued parametrization by taking a quotient of two holomorphic
functions. In both cases, the argument principle proves that the map has
degree one.

For disconnected boundary we attach a Euclidean disk along one boundary
component by smooth conformal welding. Two-dimensional Dirichlet energy
is unchanged on each piece by conformal rescaling. Minimizing over the
seam trace gives a Schur complement of the original Dirichlet-to-Neumann
map and the disk operator. The boundary arclength identification makes
the latter the square root of the seam Laplacian. This proves that
commutation survives capping, including at the constant mode. Iterating
reduces to connected boundary without changing genus.

\section{2. Holomorphic traces and the period
defect}\label{holomorphic-traces-and-the-period-defect}

We use the complex structure determined by the orientation and conformal
class of \(g\). A holomorphic function on \(M\) will mean one
holomorphic in the interior and smooth up to the boundary. Orient
\(\partial M\) positively, so that \((\nu,\tau)\) is a positively
oriented orthonormal frame when \(\tau\) is the positive unit tangent.
Set \(T=\partial_\tau\) on each component. Our Hodge-star convention is
\(*dx=dy\) in an oriented isothermal coordinate \(x+iy\).

For complex-valued traces we extend \(A\) complex linearly.
Green\textquotesingle s identity gives

\[
\int_{\partial M}\overline f\,Af\,ds
=\int_M\bigl(|du|^2+|dv|^2\bigr)\,dV_g,
\qquad u+iv=u_f.
\tag{2.1}
\]

In particular, \(A\) is self-adjoint and nonnegative on
\(L^2(\partial M,ds)\); its kernel consists of the global constants,
since \(M\) is connected. Smooth elliptic boundary regularity is
understood throughout.

\Needspace{6\baselineskip}\textbf{Lemma 2.1 (Holomorphic trace criterion).} A function
\(f\in C^\infty(\partial M;\mathbb C)\) is the trace of a holomorphic
function on \(M\) if and only if

\[
Af=-iTf.
\tag{2.2}
\]

\emph{Proof.} If \(U=u+iv\) is holomorphic, then \(dv=*du\). At the
boundary this gives \(Tv=Au\) and \(Av=-Tu\), proving (2.2).

Conversely, extend \(f\) harmonically to \(U=u+iv\). By (2.1), (2.2),
and Stokes\textquotesingle{} theorem,

\[
\begin{aligned}
\int_M(|du|^2+|dv|^2)\,dV_g
&=-i\int_{\partial M}\overline f\,df\\
&=\int_{\partial M}(u\,dv-v\,du)
=2\int_M du\wedge dv.
\end{aligned}
\tag{2.3}
\]

The imaginary term is \(-i\int_{\partial M}d(u^2+v^2)/2=0\). Pointwise,

\[
|dv-*du|^2\,dV_g
=(|du|^2+|dv|^2)\,dV_g-2du\wedge dv.
\tag{2.4}
\]

Its integral is zero by (2.3), hence \(dv=*du\). Thus \(U\) is
holomorphic. \(\square\)

Let \(\mathcal K\) be the space of these boundary traces. It contains
the constants and is closed under multiplication. Lemma 2.1 identifies
it as the kernel of \(A+iT\) on smooth complex functions. This standard
trace-algebra viewpoint is also used in the inverse problem; our use
below concerns its additional Fourier invariance under (1.2).

From now through Section 3, assume that \(\partial M\) is connected. A
constant rescaling of \(g\) preserves (1.2), since \(A\) and
\(\Delta_\partial\) scale with degrees \(-1\) and \(-2\). Normalize its
length to \(2\pi\). Positively oriented arclength is then
\(s\in\mathbb R/2\pi\mathbb Z\), and \(T=d/ds\).

Let \(C^\infty_0(\partial M;\mathbb R)\) denote the real smooth
functions with zero boundary mean. The restriction of \(A\) to this
space has a smooth inverse: the Neumann problem is solvable precisely
for mean-zero normal data, and its additive constant is fixed by the
boundary mean. Define

\[
H=A^{-1}T:C^\infty_0(\partial M;\mathbb R)
\longrightarrow C^\infty_0(\partial M;\mathbb R).
\tag{2.5}
\]

\Needspace{6\baselineskip}\textbf{Lemma 2.2 (Finite period defect).} If \(M\) has genus \(\gamma\)
and one boundary component, then \(H^2+I\) has rank at most \(2\gamma\)
on \(C^\infty_0(\partial M;\mathbb R)\).

\emph{Proof.} If \(u\) is the harmonic extension of \(f\), then \(*du\)
is closed. Consider the linear period map

\[
\mathcal P:f\longmapsto[*du]\in H^1_{\mathrm{dR}}(M;\mathbb R).
\tag{2.6}
\]

This cohomology space has dimension \(2\gamma\). If \(\mathcal P f=0\),
there is a smooth real function \(v\) with \(dv=*du\). It is harmonic,
and we choose its additive constant so that its boundary trace has mean
zero. The boundary Cauchy-Riemann identities are

\[
A(v|_{\partial M})=-Tf,\qquad T(v|_{\partial M})=Af.
\tag{2.7}
\]

Consequently \(Hf=-v|_{\partial M}\) and

\[
H^2f=-A^{-1}T(v|_{\partial M})=-f.
\tag{2.8}
\]

Thus \(H^2+I\) vanishes on \(\ker\mathcal P\) and factors through a
vector space of dimension at most \(2\gamma\). This proves the rank
bound. No existence of a harmonic conjugate is asserted outside
\(\ker\mathcal P\). \(\square\)

The exact topological rank formula for a related boundary Hilbert
transform is stronger than Lemma 2.2; see {[}BK, Section 2, formulas
(12)-(16){]} and its attribution to {[}B, Lemma 1{]}. Only the
elementary upper bound just proved is needed here.

\section{3. Connected boundary forces genus
zero}\label{connected-boundary-forces-genus-zero}

Assume (1.2), and keep the normalization \(|\partial M|=2\pi\). For
\(n\geq1\), let

\[
E_n=\operatorname{span}_{\mathbb R}\{\cos(ns),\sin(ns)\}.
\tag{3.1}
\]

These are the positive eigenspaces of \(\Delta_\partial=-T^2\). Exact
commutation implies that \(A\) preserves each \(E_n\). Self-adjointness
then implies that its orthogonal Fourier projections commute with \(A\)
on smooth functions. The operators \(A^{-1}\), \(T\), and \(H\) also
preserve \(E_n\).

\Needspace{6\baselineskip}\textbf{Lemma 3.1 (High Fourier planes).} There exists \(N\geq1\) such
that for every \(n\geq N\) there is a unique \(c_n\in\mathbb C\),
\(|c_n|<1\), for which

\[
f_n(s)=e^{ins}+c_ne^{-ins}\in\mathcal K.
\tag{3.2}
\]

\emph{Proof.} In the real orthonormal cosine-sine basis of \(E_n\),
write \(A_n\) for the positive symmetric matrix of \(A\) and
\(J=\left(\begin{smallmatrix}0&1\\-1&0\end{smallmatrix}\right)\). The
matrix of \(T\) is \(nJ\). For any positive symmetric two-by-two matrix
\(S\), direct multiplication gives

\[
(S^{-1}J)^2=-\frac{1}{\det S}I.
\tag{3.3}
\]

It follows that

\[
(H^2+I)|_{E_n}
=\left(1-\frac{n^2}{\det A_n}\right)I.
\tag{3.4}
\]

If the scalar in (3.4) were nonzero for infinitely many \(n\), the range
of \(H^2+I\) would contain infinitely many independent Fourier planes.
This contradicts Lemma 2.2. Hence

\[
\det A_n=n^2\qquad(n\geq N)
\tag{3.5}
\]

for some \(N\).

Set \(z=e^{is}\). In the complex orthonormal basis proportional to
\((z^n,z^{-n})\), reality and self-adjointness give

\[
A_n=\begin{pmatrix}a&b\\\overline b&a\end{pmatrix},\qquad a>0,
\qquad -iT=\begin{pmatrix}n&0\\0&-n\end{pmatrix}.
\tag{3.6}
\]

By (3.5), \(\det(A_n+iT)=a^2-|b|^2-n^2=0\). Its kernel is
one-dimensional, because its lower diagonal entry is \(a+n>0\). Lemma
2.1 identifies this kernel with \(\mathcal K\cap E_n^{\mathbb C}\).

For any nonzero trace \(f=\alpha z^n+\beta z^{-n}\) in this kernel,
strict positivity of (2.1) and (2.2) imply

\[
0<\int_{\partial M}\overline f\,Af\,ds
=2\pi n\bigl(|\alpha|^2-|\beta|^2\bigr).
\tag{3.7}
\]

The trace is nonconstant, so its harmonic extension has strictly
positive energy. Thus \(\alpha\ne0\) and \(|\beta/\alpha|<1\).
Normalization proves existence and uniqueness in (3.2). \(\square\)

Write \(P_n\) for the orthogonal projection onto \(E_n^{\mathbb C}\).
Since both \(A\) and \(T\) commute with \(P_n\), Lemma 2.1 gives

\[
P_n\mathcal K\subset\mathcal K\qquad(n\geq1).
\tag{3.8}
\]

\Needspace{6\baselineskip}\textbf{Proposition 3.2 (A degree-one boundary trace).} There is a
holomorphic function \(F\) on \(M\) whose boundary trace has the form

\[
F|_{\partial M}=z+c z^{-1},\qquad |c|<1.
\tag{3.9}
\]

\emph{Proof.} First suppose that \(c_n\ne0\) for some \(n\geq N\). The
product \(f_nf_{n+1}\) belongs to \(\mathcal K\). Its frequency-one
projection is

\[
P_1(f_nf_{n+1})=c_nz+c_{n+1}z^{-1}.
\tag{3.10}
\]

This is a nonzero holomorphic trace by (3.8). Applying (3.7) at
frequency one shows that \(|c_{n+1}|<|c_n|\). Dividing (3.10) by \(c_n\)
proves (3.9).

Suppose instead that \(c_n=0\) for all \(n\geq N\). Fix such an \(n\).
Let \(F_n,F_{n+1}\) be the holomorphic extensions of \(z^n,z^{n+1}\).
The holomorphic function \(F_n^{n+1}-F_{n+1}^n\) has zero boundary
trace, so the maximum principle gives

\[
F_n^{n+1}=F_{n+1}^n\quad\text{on }M.
\tag{3.11}
\]

The quotient \(G=F_{n+1}/F_n\) is meromorphic in the interior, is smooth
near the boundary since \(F_n\) is nonzero there, and satisfies
\(G^n=F_n\) wherever it is initially defined. A pole of \(G\) would give
a pole of \(G^n\), which is impossible. Its apparent singularities are
therefore removable. Thus \(G\) is holomorphic on \(M\), smooth to the
boundary, and has boundary trace \(z\). Take \(F=G\) and \(c=0\).
\(\square\)

\Needspace{6\baselineskip}\textbf{Proposition 3.3 (Connected-boundary case of Theorem 1.1).} If
\(\partial M\) is connected and (1.2) holds, then \(M\) has genus zero.

\emph{Proof.} Take \(F\) from Proposition 3.2. The real-linear map
\(w\mapsto w+c\overline w\) has determinant \(1-|c|^2>0\). Hence (3.9)
maps the oriented boundary diffeomorphically onto a positively oriented
nondegenerate ellipse \(\Gamma\). Let \(\Omega\) be the bounded
component of \(\mathbb C\setminus\Gamma\).

For \(a\notin\Gamma\), the argument principle on the bordered Riemann
surface gives

\[
\sum_{p\in F^{-1}(a)}\operatorname{ord}_p(F-a)
=\frac{1}{2\pi i}\int_{\partial M}\frac{dF}{F-a}
=\begin{cases}1,&a\in\Omega,\\0,&a\notin\overline\Omega.\end{cases}
\tag{3.12}
\]

There are finitely many zeros in (3.12), since there are none near the
boundary. For completeness, the first equality follows by deleting small
disjoint coordinate disks around those zeros and applying
Stokes\textquotesingle{} theorem to the closed form \(dF/(F-a)\) on the
remaining surface. Each small circle contributes its positive zero
multiplicity.

No interior point can map to \(\Gamma\): by the open mapping theorem its
image neighborhood would contain points outside \(\overline\Omega\),
contradicting the second case of (3.12). Therefore \(F\) maps the
interior bijectively onto \(\Omega\), and every local multiplicity is
one. It is a biholomorphism there. Its boundary map is also bijective,
so \(F:M\to\overline\Omega\) is a homeomorphism by compactness. In
particular \(M\) is a topological disk and has genus zero. \(\square\)

\section{4. Conformal capping preserves
commutation}\label{conformal-capping-preserves-commutation}

We now allow several boundary components. The following construction
concerns Dirichlet energy and keeps the line element on every uncapped
component exactly fixed.

\Needspace{6\baselineskip}\textbf{Lemma 4.1 (Smooth conformal cap).} Suppose \(M\) has at least
two boundary components, and designate one of them by \(Y\). Let \(D\)
be a Euclidean disk of perimeter \(|Y|_g\). Identify \(\partial D\) with
\(Y\) by an orientation-reversing arclength isometry. The resulting
topological surface admits a smooth conformal structure and a smooth
compatible metric \(\widehat g\) with these properties:

\begin{enumerate}
\def\labelenumi{\arabic{enumi}.}
\tightlist
\item
  The inclusions of \(M\) and \(D\) are conformal on their interiors and
  smooth with nondegenerate derivatives up to the seam.
\item
  The line element of \(\widehat g\) on \(X=\partial M\setminus Y\)
  equals the original line element of \(g\).
\item
  If \(U\in H^1(\widehat M)\) has restrictions \(u\) and \(v\) to the
  two pieces, then
\end{enumerate}

\[
\int_{\widehat M}|dU|_{\widehat g}^2\,dV_{\widehat g}
=\int_M|du|_g^2\,dV_g+\int_D|dv|^2\,dx.
\tag{4.1}
\]

The surface \(\widehat M\) is connected, has one fewer boundary
component, and has the same genus as \(M\).

\emph{Proof.} Smooth conformal sewing along an orientation-reversing
boundary diffeomorphism is standard; a general bordered-surface
statement is {[}T, Theorem 2.33{]}. The precise local smoothness can
also be obtained from disk welding {[}HT, Theorem 3.1{]} and
disk-annulus welding {[}HT, Proposition 3.6{]}. Here is the reduction to
that local form.

Choose a smooth annular collar of \(Y\) in \(M\). Its conformal
uniformization by a round annulus is smooth with nonzero derivative at
both boundary circles. In these coordinates the seam identification is a
smooth circle diffeomorphism. Welding the disk to the collar produces a
disk with conformal inclusions smooth up to a smooth seam. The conformal
structure on the open interior of the original collar is unchanged. It
therefore glues to the remainder of \(M\) along an open collar overlap.
This supplies the asserted smooth conformal structure on \(\widehat M\).

Choose any smooth metric \(h\) compatible with this structure. Its
restriction to \(M\) is conformal to \(g\), and its restriction to \(D\)
is conformal to the Euclidean metric. On each component of \(X\),
prescribe a positive smooth scalar factor so that a conformal rescaling
of \(h\) has the original line element. Extend this factor smoothly and
positively to \(\widehat M\), using boundary collars, and denote the
rescaled metric by \(\widehat g\).

In dimension two,
\(|du|_{e^{2\varphi}g}^2dV_{e^{2\varphi}g}=|du|_g^2dV_g\). Applying this
on both pieces proves (4.1). Matching \(H^{1/2}\) traces across the
smooth seam are precisely the compatibility condition for two piecewise
\(H^1\) functions to define an \(H^1\) function on \(\widehat M\). This
follows in a smooth seam chart from the usual trace and gluing property
for a hypersurface; a jump is the only additional distributional
derivative. No matching of normal metric derivatives across the seam is
required. Finally, attaching a disk increases Euler characteristic by
one and decreases the boundary count by one, so the genus is unchanged.
\(\square\)

The arclength condition determines the disk operator transported to the
seam:

\[
B=\Delta_Y^{1/2}.
\tag{4.2}
\]

Indeed, on a disk of radius \(r=|Y|_g/(2\pi)\), the modes
\(e^{in\theta}\) have Dirichlet-to-Neumann eigenvalues \(|n|/r\) and
boundary Laplace eigenvalues \(n^2/r^2\). Reversing orientation leaves
both operators unchanged.

Relative to \(L^2(X)\oplus L^2(Y)\) with the original line elements,
write the original operator in blocks as

\[
A=\begin{pmatrix}A_{XX}&A_{XY}\\A_{YX}&A_{YY}\end{pmatrix}.
\tag{4.3}
\]

\Needspace{6\baselineskip}\textbf{Proposition 4.2 (Capping and commutation).} If \(A\) commutes
with \(\Delta_\partial\), then the Dirichlet-to-Neumann map
\(\widehat A\) of \((\widehat M,\widehat g)\) in Lemma 4.1 commutes with
the unchanged boundary Laplacian \(\Delta_X\).

\emph{Proof.} Let \(d_M\) and \(d_{\widehat M}\) denote the minimized
Dirichlet-energy forms for prescribed boundary values. Formula (4.1) and
the trace gluing property give, for \(f\in H^{1/2}(X)\),

\[
d_{\widehat M}[f]
=\inf_{h\in H^{1/2}(Y)}
\left\{d_M[(f,h)]+\langle Bh,h\rangle\right\}.
\tag{4.4}
\]

The brackets in (4.4) denote Sobolev duality, agreeing with the \(L^2\)
energy pairing on smooth real functions; complex forms follow by
polarization. Minimization over the interior of each piece first
produces its harmonic extension, which justifies both directions of
(4.4).

The seam operator

\[
C=A_{YY}+B:H^{1/2}(Y)\longrightarrow H^{-1/2}(Y)
\tag{4.5}
\]

is coercive. In fact, if \(u_h\) has boundary trace \((0,h)\) on
\(X\sqcup Y\), then the Poincare inequality with zero trace on the
nonempty boundary portion \(X\), followed by the trace theorem, yields

\[
\|h\|_{H^{1/2}(Y)}^2
\leq C_0\|u_h\|_{H^1(M)}^2
\leq C_1\int_M|du_h|^2\,dV_g
=C_1\langle A_{YY}h,h\rangle.
\tag{4.6}
\]

Adding the nonnegative disk energy preserves coercivity. The bounded
form of \(C\) is therefore invertible by the variational existence
theorem. This includes constants on \(Y\): although \(B\) annihilates
constants, the extension with zero trace on \(X\) supplies the positive
term in (4.6).

The unique minimizing trace in (4.4) satisfies

\[
Ch=-A_{YX}f.
\tag{4.7}
\]

Polarization, or differentiation of the minimum energy with respect to
\(f\), gives the exact Schur complement

\[
\widehat A=A_{XX}-A_{XY}C^{-1}A_{YX}.
\tag{4.8}
\]

The boundary measure used here is the original measure on \(X\), which
equals the one for \(\widehat g\) by Lemma 4.1. Thus (4.8) is the
geometric Dirichlet-to-Neumann map, with no additional boundary
multiplier.

It remains to check commutation without an unjustified manipulation of
unbounded operators. The original commutation relation gives the block
identities

\[
\Delta_iA_{ij}=A_{ij}\Delta_j,
\qquad i,j\in\{X,Y\},
\tag{4.9}
\]

on smooth functions. The diagonal block \(A_{YY}\) is self-adjoint and
preserves every \(\Delta_Y\) eigenspace, so it commutes with its
orthogonal spectral projections. The same is true of \(B\) and hence of
the form operator \(C\). These projection identities extend to the
indicated Sobolev spaces by density. Its inverse also commutes with the
projections: apply a projection to \(Ch=k\) and use uniqueness of the
solution.

Now let \(f\) belong to the \(\lambda\) eigenspace of \(\Delta_X\). By
(4.9), \(A_{YX}f\) belongs to the \(\lambda\) eigenspace of
\(\Delta_Y\), or is zero if that eigenvalue is absent there. Equation
(4.7) therefore puts \(h\) in the same eigenspace. Equations (4.8)-(4.9)
imply that \(\widehat Af\) is in the \(\lambda\) eigenspace of
\(\Delta_X\). This also holds for \(\lambda=0\) by the invertibility
already proved. Hence \(\widehat A\) preserves every \(\Delta_X\)
eigenspace. Smooth Fourier expansion on the finite union of boundary
circles, together with continuity of both operators on smooth functions,
proves their commutation. \(\square\)

\emph{Proof of Theorem 1.1.} If the boundary is connected, use
Proposition 3.3. Otherwise cap one boundary component using Lemma 4.1
and Proposition 4.2. The resulting surface has the same genus, one fewer
boundary component, and commuting boundary operators. Repeat until the
boundary is connected. Proposition 3.3 makes the final genus zero, hence
the original genus was zero. \(\square\)

\section{5. Scope and further
questions}\label{scope-and-further-questions}

The proof uses exact commutation twice: first to reduce the
finite-period defect to scalar defects on Fourier planes, and then to
preserve the holomorphic-trace algebra under projection. A small
commutator does not imply either exact determinant identities or exact
multiplicative Fourier relations. A quantitative genus obstruction would
require additional geometric control and is not a consequence of this
argument. The stability results in {[}S2{]} concern other settings and
do not supply that conclusion for arbitrary bordered surfaces.

Proposition 4.2 may be useful separately when studying boundary
symmetries compatible with conformal sewing. Its hypotheses specify the
metric on the uncapped boundary and the arclength identification at the
seam; omitting either can change the operators in (4.8). The present
theorem concerns orientable surfaces with smooth metrics and smooth
boundary. We make no classification claim for nonorientable surfaces or
for rough boundaries.

\textbf{Research and preparation statement.} This manuscript was
developed and written with AI assistance, including its proof
development, source comparison, and internal audit. It has not received
independent expert review. The accompanying source and review records
distinguish prior results, the proposed contribution, and the limits of
the priority search. No numerical computation is used as evidence for
the theorem.

\sloppy\section{References}\label{references}

{[}B{]} M. I. Belishev, \emph{The Calderon problem for two-dimensional
manifolds by the BC-method}, SIAM Journal on Mathematical Analysis
\textbf{35} (2003), no. 1, 172-182.
\href{https://doi.org/10.1137/S0036141002413919}{DOI:
10.1137/S0036141002413919}.

{[}BK{]} M. I. Belishev and D. V. Korikov, \emph{On stability of
determination of Riemann surface from its DN-map},
\href{https://arxiv.org/abs/2112.14816v2}{arXiv:2112.14816v2}, 2022.
Section 2, especially formulas (12)-(16).

{[}HT{]} A. G. Henriques and J. E. Tener, \emph{The Segal-Neretin
semigroup of annuli},
\href{https://arxiv.org/abs/2410.05929v1}{arXiv:2410.05929v1}, 2024.
Theorem 3.1 and Proposition 3.6.

{[}S{]} R. Speciel, \emph{Surfaces with commuting boundary Laplacian and
Dirichlet-to-Neumann map}, Journal of Spectral Theory \textbf{16}
(2026), 1281-1292. \href{https://doi.org/10.4171/JST/617}{DOI:
10.4171/JST/617}. Final journal numbering is used throughout.

{[}S2{]} R. Speciel, \emph{Stability for commutativity properties of the
Dirichlet-to-Neumann map}, International Mathematics Research Notices
\textbf{2026} (2026), no. 14, rnag155.
\href{https://doi.org/10.1093/imrn/rnag155}{DOI: 10.1093/imrn/rnag155}.
Preprint titled \emph{Stability Estimates for Commutativity Properties
of the Dirichlet-to-Neumann Map},
\href{https://arxiv.org/abs/2510.08822v1}{arXiv:2510.08822v1}.

{[}T{]} J. E. Tener, \emph{Construction of the unitary free fermion
Segal CFT}, Communications in Mathematical Physics \textbf{355} (2017),
463-518. \href{https://doi.org/10.1007/s00220-017-2959-x}{DOI:
10.1007/s00220-017-2959-x}. The sewing statement used here is Theorem
2.33 in the inspected author preprint,
\href{https://arxiv.org/abs/1608.02095v1}{arXiv:1608.02095v1}, 2016.

\end{document}
